\chapter{Conclusiones y Discusión Grupal}

A partir de la experiencia de laboratorio, el grupo de trabajo expone las siguientes resoluciones:

\begin{itemize}
    \item \textbf{Sentido de las pasadas:} ¿Cuál es el sentido de efectuar una pasada hacia arriba y otra hacia abajo durante la contrastación del voltímetro? ¿Sería mejor promediar varias series de pasadas (cuatro o cinco)? ¿Supondría ventaja o mejora?
    \\ \textit{[Respuesta...]}

    \item \textbf{Taller de mantenimiento:} Averigüe si en el Taller de mantenimiento del laboratorio el instrumental es sometido a contrastaciones periódicas.
    \\ \textit{[Respuesta...]}

    \item \textbf{Instrumentos de referencia en el laboratorio:} Elabore una lista de al menos tres instrumentos que existan en el laboratorio que podrían ser usados como referencia (patrón), incluyendo su especificación de exactitud.
    \\ \textit{[Respuesta...]}

    \item \textbf{Validez de la curva:} En el instrumento UNIVO contrastado, ¿Por cuánto tiempo estima que será válida la curva obtenida? ¿Cuáles serían las causas que obligarían a efectuar un nuevo procedimiento de contrastación?
    \\ \textit{[Respuesta...]}

    \item \textbf{Métodos de 4 terminales:} ¿Qué ventajas presentan los métodos de medición con 4 terminales por sobre los métodos convencionales (óhmetro de 2 puntas)?
    \\ \textit{[Respuesta...]}

    \item \textbf{Caja de resistencias de precisión:} Realice una búsqueda sobre "caja de resistencias de precisión". Busque hojas de especificaciones e interprete las correspondientes a la incertidumbre.
    \\ \textit{[Respuesta...]}

    \item \textbf{Fuente SK150 (LM317):} La fuente de corriente propuesta es simple pero efectiva. Busque la hoja de datos del Circuito integrado LM317 e investigue cómo funciona y cómo mantiene la corriente constante.
    \\ \textit{[Respuesta...]}

    \item \textbf{Propagación de incertidumbre:} La ecuación empleada asume que los errores parciales se suman (peor caso). Algunos autores sugieren que se obtiene una mejor estimación efectuando la \textit{raíz cuadrada de la suma de los cuadrados} de los errores parciales. Compare ambos enfoques.
    \\ \textit{[Respuesta...]}
\end{itemize}
